\documentclass[10pt,a4paper]{amsart}
\usepackage[margin=19mm]{geometry}
\usepackage{amsmath,amssymb,mathtools}
\usepackage[scr=rsfs,cal=euler]{mathalfa}
\usepackage{xfrac}
\usepackage[unicode,hidelinks,bookmarks=false]{hyperref}
\newcommand{\ie}{i.e.\ }
\newcommand{\cf}{cf.\ }
\newcommand{\resp}{respectively\ }
\newcommand{\Subsection}{\subsection}
\providecommand{\nobreakdash}{\nobreak\hbox{-}}
\setlength{\emergencystretch}{2em}
\multlinegap0pt
\usepackage{bm}
\usepackage{bbm}
\usepackage{dsfont}
\usepackage{xspace}
\usepackage{cancel}
\usepackage{mathtools}
\usepackage{amsthm}
\makeatletter
\@ifundefined{theorem}{\newtheorem{theorem}{Theorem}}{}
\@ifundefined{claim}{\newtheorem{claim}[theorem]{Claim}}{}
\@ifundefined{lemma}{\newtheorem{lemma}[theorem]{Lemma}}{}
\makeatother
\newcommand{\R}{\mathbb R}
\newcommand{\Z}{\mathbb Z}
\newcommand{\conv}{{\rm conv}\,}
\renewcommand{\ge}{\geqslant}
\renewcommand{\le}{\leqslant}
\title[Convex geometry and the Erd\H{o}s-Ginzburg-Ziv problem]{Convex geometry and the Erd\H{o}s-Ginzburg-Ziv problem}
\author{Dmitrii Zakharov}
\date{Source: arXiv:2002.09892v6 (CC BY 4.0)}
\begin{document}
\maketitle

\noindent\emph{Source and licence.} Convex geometry and the Erd\H{o}s-Ginzburg-Ziv problem. Version arXiv:2002.09892v6 (CC BY 4.0), distributed under the Creative Commons Attribution 4.0 International licence, as recorded in the arXiv licence metadata for this exact version and shown on the abstract page https://arxiv.org/abs/2002.09892v6. Full rights notice: licenses/arxiv-2002.09892-CC-BY-4.0.txt.\par\smallskip

\noindent\textit{Editorial scope.} Counted unit: the Lemma whose source label is \texttt{lm2} in the article (source0.tex lines 1668--1680, source label \texttt{lm2}), delivered together with the article's own complete original proof of it (source0.tex lines 1683--1761, one \texttt{proof} environment of 79 lines). No mathematics is written, repaired or re-derived: statement and proof are the article's own text line for line. Required context retained uncounted: none. Every definition, display and label of those lines is retained unchanged. Omitted from this package and not redistributed: 1--1667, 1681--1682, 1762--1963. Nothing inside a retained excerpt is omitted or altered: every retained body is delivered exactly as the article's lines are written, so no omission receipt of any kind accompanies this package. Citations: the retained excerpts carry no citation command at all, so no bibliography entry of the article is transcribed and no reference is redistributed. Reference adaptation: none was needed and none was made; every reference inside the delivered unit resolves inside it, so no unresolved reference or citation is produced. Printed slips kept literal: no printed word, symbol, hypothesis or number of the counted statement or of its proof was altered. Layout and class: the article's own class is \texttt{daj} with packages including amsfonts, amsmath, amsthm, amssymb, indentfirst, xy, adjustbox, subcaption, wrapfig, thmtools, thm-restate, babel; the wrapper is the lane's pinned \texttt{amsart} editorial wrapper at 10pt on A4, which restates the theorem-like environments the retained text needs and adds the article's own macro definitions that the retained text needs (\texttt{\textbackslash R}, \texttt{\textbackslash Z}, \texttt{\textbackslash conv}, \texttt{\textbackslash ge}, \texttt{\textbackslash le}), with extra wrapper packages (bm, bbm, dsfont, xspace, cancel, mathtools). The delivered numbering is the unit's own. Assets: no retained span contains a figure environment, a graphics-inclusion command, a PDF-page inclusion, a table or a TikZ picture; no image file is copied into this package, the package declares no asset dependency and no image file is redistributed with it. Licence: version arXiv:2002.09892v6 of this article is distributed under the Creative Commons Attribution 4.0 International licence, as recorded in the arXiv licence metadata for this exact version and shown on the abstract page https://arxiv.org/abs/2002.09892v6. A full rights notice is delivered at licenses/arxiv-2002.09892-CC-BY-4.0.txt. Characters: every retained excerpt is pure ASCII, so the delivered engine, XeLaTeX with its default Latin Modern text font, reports no missing character.
\begin{lemma}\label{lm2}
Let $\theta > 0$, let $w: \R^d \rightarrow \R_{\ge 0}$ be a function with finite support $S$. 
Let $\Lambda$ be the minimal lattice containing $S$ and $c \in \Lambda \cap \,{\rm int}(\conv{S})$ be a $\theta$-central point for $w$.

Then for every $\varepsilon > 0$ and all $n > n_0(\varepsilon, w, c)$ there are non-negative integer coefficients $\alpha_q$ for $q \in S$ and $\mu = \mu(\varepsilon, w, c) > 0$ such that:
\begin{align}\label{eqlm}
    \sum_{q \in S} \alpha_q = n,~~    \sum_{q \in S} \alpha_q q = n c, 
\end{align}
and for every $q\in S$ we have 
\begin{align}\label{eqlm2}
    \mu n \le \alpha_q \le  (1 + \varepsilon) \frac{n w(q)}{\theta w(S)}.
\end{align}
\end{lemma}

\begin{proof}
Without loss of generality, we may assume that $c = 0$, the set $S$ spans $\R^d$, $\Lambda = \Z^d$ and $w(S) = \sum_{q \in S} w(q) = 1$.

\begin{claim}\label{real}
There are rational coefficients $\beta_q$ such that:
\begin{align*}
    \sum_{q\in S} \beta_q q = 0,~~ \sum_{q \in S} \beta_q = 1, 
\end{align*}
and $\beta_q \in (0, \theta^{-1}w(q))$ for every $q\in S$.
\end{claim}
\begin{proof}

Note that it is enough to find {\em real} coefficients $\beta_q$ with properties described in the claim. The existence of rational coefficients would then follow automatically.

We denote by $\R^S$ the space of all functions $\xi: S \rightarrow \R$. This space is equipped with the natural scalar product $\xi \cdot \eta = \sum_{q \in S} \xi(q) \eta(q)$. In what follows we identify $\R^S$ with the dual space $(\R^S)^*$ via this scalar product.

Let $H \subset \R^S$ be the set of vectors $(c_q)_{q \in S}$ such that $\sum_{q \in S} c_q q = 0$. 
Let $\Omega \subset \R^S$ be the set of all functions $v: S \rightarrow \R$ such that 
$$
0 \le v(q) \le \theta^{-1}w(q) \sum_{q' \in S} v(q'),
$$
for every $q \in S$. Note that if the intersection $H \cap \operatorname{int}(\Omega)$ is non-empty, then we are done: take a vector $v \in H \cap \operatorname{int}(\Omega)$ and define $\beta_q = \frac{v(q)}{\sum_{q' \in S} v(q')}$.

Let us assume that $H \cap \operatorname{int}(\Omega) = \emptyset$ and arrive at a contradiction. Since $H$ is a vector subspace and $\operatorname{int}(\Omega)$ is an open convex set, there exists a function $\xi \in \R^S$ such that  
$$
\xi(H) = 0 ~~ {\text{and}} ~~\xi(\Omega) \ge 0.
$$
The first condition can be reformulated as $\xi \in H^\bot$.
Note that the space $H^\bot$ is isomorphic to $\R^d$: given a function $\zeta \in H^\bot$ we define a linear function $\tilde \zeta$ on $\R^d$ by setting $\tilde \zeta(q) = \zeta(q)$ for $q \in S$ and extending $\tilde \zeta$ by linearity. The conditions that $S$ spans $\R^d$ and the linear equations defining $H^{\bot}$ imply that this definition is correct. 
Let $\tilde \xi$ be the linear function on $\R^d$ corresponding to $\xi$.

Let $\varepsilon_q$ be the element of the standard basis of $\R^S$ corresponding to $q \in S$ and denote $\sigma = \sum_{q \in S}\varepsilon_q$. The set $\Omega$ is defined as the set of vectors $v \in \R^S$ such that
\begin{equation} \label{defw}
\varepsilon_q \cdot v \ge 0~~ \text{and} ~~(w(q)\sigma - \theta \varepsilon_q) \cdot v \ge 0,
\end{equation}

for all $q \in S$. By duality, the condition $\xi(\Omega) \ge 0$ is a non-negative linear combination of inequalities (\ref{defw}).
Indeed, if not, then $\xi$ can be separated by a hyperplane from functions (\ref{defw}) in the space of all linear functions on $\R^S$. But this hyperplane will correspond to a point in $\Omega$ on which the value of $\xi$ is negative.
Thus, there are nonnegative real coefficients $a_q, b_q \ge 0$ such that
\begin{equation}\label{eqxi1}
    \xi = \sum_{q \in S} a_q \varepsilon_q + b_q (w(q)\sigma - \theta \varepsilon_q) = \sum_{q \in S} (a_q - \theta b_q)\varepsilon_q + \left( \sum_{q \in S} b_q w(q) \right) \sigma.
\end{equation}

Let $I \subset S$ be the set of $q \in S$ such that $\xi \cdot \varepsilon_q \le 0$. Since $c= 0$ is a $\theta$-central point for $w$ and $\xi\cdot \varepsilon_q = \tilde \xi(q)$ for all $q \in S$, we have 
\begin{equation}\label{omegaq}
\sum_{q \in I} w(q) \ge \theta. 
\end{equation}
On the other hand, for every $q \in I$ by (\ref{eqxi1}) we have
\begin{equation}\label{xieq}
\xi(q) = (a_q - \theta b_q) + \left(\sum_{q' \in S} b_{q'} w(q') \right) \le 0,
\end{equation}
hence, by discarding the non-negative term $a_q$ from (\ref{xieq}) we get
$$
\theta b_q \ge \sum_{q' \in S} b_{q'} w(q').
$$
Summing this over $q \in I$ with weights $w(q)>0$ we obtain:
$$
\theta \sum_{q \in I} b_q w(q) \ge \left( \sum_{q \in I} w(q) \right) \left( \sum_{q \in S} b_{q} w(q)  \right) \stackrel{(\ref{omegaq})}{\ge} \theta \left( \sum_{q \in S} b_{q} w(q)  \right).
$$
The sum on the left hand side is a subsum of the right hand side. Since $\theta > 0$ and $w(q) > 0$ for all $q$, it follows that an equality is attained in (\ref{xieq}) for all $q \in I$. This, however, means that $\xi(q) \ge 0$ for every $q\in S$ and the point $c = 0$ lies on the boundary of $\conv(S)$ which contradicts our assumption.
We conclude that there cannot be such a function $\xi$ and hence $H \cap {\rm int}(\Omega) \neq \emptyset$, as desired.
\end{proof}

Take some rational coefficients $\beta_q$ provided by Claim \ref{real}; note that we can define them as functions of $w$ and $c$. Let $m$ be the least common multiple of denominators of $\beta_q$.
Since $c=0$ belongs to the minimal lattice of $S$ there is an integer vector $\delta \in \Z^S$ such that $\sum_{q \in S} \delta_q q = 0$ and $\sum_{q \in S} \delta_q = 1$. Let $C = \max_{q \in S} |\delta_q|$. 

Let us define the function $n_0 = n_0(\varepsilon, w, c)$ by
$$
n_0 = 2Cm^2 + \varepsilon^{-1} Cm\theta \max_{q\in S} w(q)^{-1}, 
$$
(note that $w(q) > 0$ for every $q \in S$ by assumption). Now consider an arbitrary $n > n_0$. Write $n = am+r$ where $0 \le r < m$ and define the coefficients by $\alpha_q = a m \beta_q + r \delta_q$; note that $\alpha_q$ is an integer.
Let us check that all required conditions are satisfied:
\begin{align*}
    \sum_{q \in S} \alpha_q q &= \sum_{q \in S} am\beta_q q + r \delta_q q = 0, \\
    \sum_{q \in S} \alpha_q &= am+r = n, \\
    \alpha_q = am\beta_q + r \delta_q &\le am \theta^{-1} w(q) + rC \le n \theta^{-1} w(q) (1 + mC n^{-1} \theta w(q)^{-1}) < n\theta^{-1} w(q) (1 + \varepsilon),
\end{align*}
A similar computation gives $\alpha_q > \mu n$ for some $\mu > 0$ not depending on $n$. Lemma \ref{lm2} is proved.
\end{proof}

\end{document}
